% GENERATED FILE. Edit canonical lessons, then run npm run generate:books.
% canonical-source-hash: fnv1a64:2c719ed9aef9b717
% canonical-lessons: TE-C199-dancu, TE-C199-nimpu, TE-C199-arabettu, TE-C199-madatapettu, TE-C199-sardu

\chapter{To pound, To fill, To put out to dry, To fold, To arrange}
\label{ch:te-to-pound-to-fill-to-put-out-to-dry-to-fold-to-arrange}
\hlchaptermodality{199}

\begin{chapteropening}
\textbf{By the end of this chapter:} \emph{I can say to pound, to fill, to put out to dry, to fold and to arrange.}

Complete the last lesson of chapter 199: I can say to pound, to fill, to put out to dry, to fold and to arrange.
\end{chapteropening}

\section[\texorpdfstring{\te{దంచు} (dañcu)}{దంచు (dañcu)}]{\te{దంచు} (dañcu) — to pound (in a mortar)}

\label{lesson:TE-C199-dancu}



\begin{quote}
Before the new one: say the Telugu for a uniform, then the Telugu for staff.
\end{quote}

\subsection*{You'll want to know: \te{దంచు}}
\textbf{\te{దంచు}} — \emph{dañcu} — \textquotedblleft{}to pound (in a mortar)\textquotedblright{}.

\begin{grammarlens}[title={What you've built}]
One more word for this chapter.
\end{grammarlens}

\subsection*{Your turn}
\begin{itemize}
\raggedright
  \item \emph{Say it:} \emph{dañcu}
  \item \emph{Say it:} \emph{dañcu}, once more
  \item \emph{Say it:} say \emph{dañcu}
  \item \emph{Recall:} say the Telugu for a uniform, then the Telugu for staff, then say \emph{dañcu} again
\end{itemize}

\begin{tcolorbox}[breakable,colback=teal!4,colframe=teal!35!black,title={Before you move on}]
What does \te{దంచు} mean? (\textquotedblleft{}To pound (in a mortar)\textquotedblright{}.) Say it once more.
\end{tcolorbox}
\section[\texorpdfstring{\te{నింపు} (nimpu)}{నింపు (nimpu)}]{\te{నింపు} (nimpu) — to fill}

\label{lesson:TE-C199-nimpu}



\begin{quote}
Before the new one: say the Telugu for staff, then the Telugu for to pound.
\end{quote}

\subsection*{You'll want to know: \te{నింపు}}
\textbf{\te{నింపు}} — \emph{nimpu} — \textquotedblleft{}to fill\textquotedblright{}.

\begin{grammarlens}[title={What you've built}]
One more word for this chapter.
\end{grammarlens}

\subsection*{Your turn}
\begin{itemize}
\raggedright
  \item \emph{Say it:} \emph{nimpu}
  \item \emph{Say it:} \emph{nimpu}, once more
  \item \emph{Say it:} say \emph{nimpu}
  \item \emph{Recall:} say the Telugu for staff, then the Telugu for to pound, then say \emph{nimpu} again
\end{itemize}

\begin{tcolorbox}[breakable,colback=teal!4,colframe=teal!35!black,title={Before you move on}]
What does \te{నింపు} mean? (\textquotedblleft{}To fill\textquotedblright{}.) Say it once more.
\end{tcolorbox}
\section[\texorpdfstring{\te{ఆరబెట్టు} (ārabeṭṭu)}{ఆరబెట్టు (ārabeṭṭu)}]{\te{ఆరబెట్టు} (ārabeṭṭu) — to put out to dry}

\label{lesson:TE-C199-arabettu}



\begin{quote}
Before the new one: say the Telugu for to pound, then the Telugu for to fill.
\end{quote}

\subsection*{You'll want to know: \te{ఆరబెట్టు}}
\textbf{\te{ఆరబెట్టు}} — \emph{ārabeṭṭu} — \textquotedblleft{}to put out to dry\textquotedblright{}.

\begin{grammarlens}[title={What you've built}]
One more word for this chapter.
\end{grammarlens}

\subsection*{Your turn}
\begin{itemize}
\raggedright
  \item \emph{Say it:} \emph{ārabeṭṭu}
  \item \emph{Say it:} \emph{ārabeṭṭu}, once more
  \item \emph{Say it:} say \emph{ārabeṭṭu}
  \item \emph{Recall:} say the Telugu for to pound, then the Telugu for to fill, then say \emph{ārabeṭṭu} again
\end{itemize}

\begin{tcolorbox}[breakable,colback=teal!4,colframe=teal!35!black,title={Before you move on}]
What does \te{ఆరబెట్టు} mean? (\textquotedblleft{}To put out to dry\textquotedblright{}.) Say it once more.
\end{tcolorbox}
\section[\texorpdfstring{\te{మడతపెట్టు} (maḍatapeṭṭu)}{మడతపెట్టు (maḍatapeṭṭu)}]{\te{మడతపెట్టు} (maḍatapeṭṭu) — to fold}

\label{lesson:TE-C199-madatapettu}



\begin{quote}
Before the new one: say the Telugu for to fill, then the Telugu for to put out to dry.
\end{quote}

\subsection*{You'll want to know: \te{మడతపెట్టు}}
\textbf{\te{మడతపెట్టు}} — \emph{maḍatapeṭṭu} — \textquotedblleft{}to fold\textquotedblright{}.

\begin{grammarlens}[title={What you've built}]
One more word for this chapter.
\end{grammarlens}

\subsection*{Your turn}
\begin{itemize}
\raggedright
  \item \emph{Say it:} \emph{maḍatapeṭṭu}
  \item \emph{Say it:} \emph{maḍatapeṭṭu}, once more
  \item \emph{Say it:} say \emph{maḍatapeṭṭu}
  \item \emph{Recall:} say the Telugu for to fill, then the Telugu for to put out to dry, then say \emph{maḍatapeṭṭu} again
\end{itemize}

\begin{tcolorbox}[breakable,colback=teal!4,colframe=teal!35!black,title={Before you move on}]
What does \te{మడతపెట్టు} mean? (\textquotedblleft{}To fold\textquotedblright{}.) Say it once more.
\end{tcolorbox}
\section[\texorpdfstring{\te{సర్దు} (sardu)}{సర్దు (sardu)}]{\te{సర్దు} (sardu) — to arrange, to tidy up}

\label{lesson:TE-C199-sardu}



\begin{quote}
Before the new one: say the Telugu for to put out to dry, then the Telugu for to fold.
\end{quote}

\subsection*{You'll want to know: \te{సర్దు}}
\textbf{\te{సర్దు}} — \emph{sardu} — \textquotedblleft{}to arrange, to tidy up\textquotedblright{}.

\begin{grammarlens}[title={What you've built}]
One more word for this chapter.
\end{grammarlens}

\subsection*{Your turn}
\begin{itemize}
\raggedright
  \item \emph{Say it:} \emph{sardu}
  \item \emph{Say it:} \emph{sardu}, once more
  \item \emph{Say it:} say \emph{sardu}
  \item \emph{Recall:} say the Telugu for to put out to dry, then the Telugu for to fold, then say \emph{sardu} again
\end{itemize}

\begin{tcolorbox}[breakable,colback=teal!4,colframe=teal!35!black,title={Before you move on}]
What does \te{సర్దు} mean? (\textquotedblleft{}To arrange, to tidy up\textquotedblright{}.) Say it once more.
\end{tcolorbox}
